\section{Related Work}


Just-in-time (JIT) defect prediction aims to estimate whether a code change is likely to introduce a defect at the moment it is submitted. Kim et al.\ introduced change classification using features derived from change size, diffusion, and file entropy~\cite{kim2008classifying}, establishing the per-commit granularity as a practical setting for early defect detection. Kamei et al.\ conducted large-scale empirical studies demonstrating that inspecting a small portion of commits can prevent a substantial fraction of defect-inducing changes~\cite{kamei2012large}, with change size and complexity metrics (diff-aware features) emerging as the strongest predictors while developer experience (diff-opaque) provided complementary signal. Their cross-project analysis~\cite{kamei2016studying} showed JIT models can transfer across projects under certain conditions, relevant to cold-start scenarios where historical data is unavailable. Deep learning approaches have been applied to learn semantic representations: Hoang et al.\ proposed DeepJIT, jointly encoding commit messages and diffs using CNNs~\cite{hoang2019deepjit}, with ablation studies showing code-level features alone achieved 85--90\% of full model performance, demonstrating that diff-aware features carry substantial predictive power independently. Pre-trained models such as CodeBERT~\cite{feng2020codebert} have been adapted to defect prediction, but these approaches use models as end-to-end classifiers requiring fine-tuning and large corpora. In contrast, we use LLMs strictly as feature extractors---leveraging their multi-language understanding to compute structured metrics, then feeding those into lightweight ML classifiers. This avoids the cost and opacity of end-to-end models while retaining language generalization. Other research investigates finer granularity: Pascarella et al.\ explored file-level JIT prediction~\cite{pascarella2019fine}, and Zeng et al.\ highlighted reproducibility issues in deep JIT approaches~\cite{zeng2021deep}.

Despite extensive work, many models depend on developer-specific metadata, detailed historical traces, or computation-heavy deep learning pipelines. Such dependencies create challenges for industrial adoption, particularly where privacy constraints and rapid onboarding of new projects limit data availability. Our work complements prior research by focusing strictly on \emph{diff-only} features, enabling a metadata-free and lightweight approach suitable for operational settings such as Prime Video's live-event streaming environment. Our findings validate that diff features alone can achieve strong performance (F1: 0.846), directly addressing cold-start scenarios common in industrial settings with frequent repository creation and strict privacy requirements.

